\chapter{Σχήμα Αρχείων Αποτελεσμάτων}
\label{app:schema}

Κάθε γραμμή του αρχείου μετρήσεων \path{results/raw/benchmarks_seed_2026_thesis.jsonl} είναι ένα αντικείμενο JSON με τα ίδια 22 πεδία:

\begin{verbatim}
case_id, profile, seed, dataset, scramble, scramble_depth,
state, known_exact_distance,
distance_kind, distance_value, distance_method,
solver, solution, solution_length, metric,
runtime_seconds, expanded_nodes, generated_nodes,
table_size_bytes, status, verified, notes
\end{verbatim}

Το σενάριο \code{scripts/verify_results.py} κάνει τρεις ελέγχους. Πρώτον, ότι όλα τα πεδία υπάρχουν. Δεύτερον, ότι κάθε γραμμή με \code{verified=true} ελέγχεται ξανά από τον ανεξάρτητο επαληθευτή. Τρίτον, ότι οι ακριβείς ρηχές λύσεις συμφωνούν με εξαντλητική αναζήτηση κατά πλάτος (\eng{breadth-first search}, BFS), όπου υπάρχει διαθέσιμη απόδειξη. Ελέγχει επίσης την ακεραιότητα των πινάκων προτύπων (\eng{pattern databases}, PDB): ότι ο γωνιακός πίνακας έχει 88\,179\,840 εγγραφές και ότι κάθε πίνακας ακμών έχει 42\,577\,920 εγγραφές ανά υποσύνολο. Επιπλέον, ότι τα δυαδικά αρχεία είναι πλήρη και τα αθροίσματα ελέγχου σωστά. Τέλος, επιβεβαιώνει τρία πράγματα για τον κύβο 2×2×2: ότι η σύνοψή του (\path{results/processed/pocket_cube_summary_seed_2026_thesis.json}) είναι πλήρης, ότι το άθροισμα της πλήρους κατανομής (\path{results/raw/pocket_cube_distribution_seed_2026_thesis.json}) ισούται με 3\,674\,160, και ότι οι αντιπροσωπευτικές του λύσεις είναι επαληθευμένες.

\section{Πεδία ταυτότητας περίπτωσης}
\label{sec:schema-identity}
Τα πεδία \code{case_id}, \code{profile}, \code{seed}, \code{dataset} και \code{scramble_depth} ταυτοποιούν την πειραματική περίπτωση. Το \code{case_id} είναι περιγραφικό και ευανάγνωστο, όπως \code{solved}, \code{shallow_3} ή \code{random_1_10}. Το \code{dataset} ομαδοποιεί τις περιπτώσεις σε τρία σύνολα: \code{A} για τη λυμένη κατάσταση, \code{B} για τις ρηχές ντετερμινιστικές περιπτώσεις και \code{C} για τις βαθύτερες τυχαίες περιπτώσεις. Με το \code{seed} και το \code{scramble} η περίπτωση μπορεί να αναπαραχθεί χωρίς εξωτερικό αρχείο εισόδου.

\section{Πεδία κατάστασης και απόστασης}
\label{sec:schema-distance}
Το πεδίο \code{state} αποθηκεύει τη συμβολοσειρά ψηφίδων (\eng{facelet string}) της αρχικής κατάστασης. Έτσι ο επαληθευτής μπορεί να ξαναφορτώσει την κατάσταση χωρίς να γνωρίζει το αρχικό ανακάτεμα (\eng{scramble}). Το πεδίο \code{known_exact_distance} δίνει την ακριβή απόσταση της περίπτωσης, όταν αυτή έχει αποδειχθεί από πριν με εξαντλητική απόδειξη (στις ρηχές περιπτώσεις, έως απόσταση 5). Αλλιώς μένει κενό. Τα πεδία \code{distance_kind}, \code{distance_value} και \code{distance_method} περιγράφουν το αποτέλεσμα του προσδιορισμού της απόστασης. Αν το \code{distance_kind} είναι \code{exact_distance}, η τιμή είναι αποδεδειγμένα ακριβής, μέσα στα όρια βάθους της αναζήτησης. Αν είναι \code{lower_bound}, η τιμή είναι μόνο κάτω φράγμα: αποκλείει μικρότερες αποστάσεις, αλλά δεν είναι η ίδια η πραγματική απόσταση. Το \code{distance_method} καταγράφει τη μέθοδο που παρήγαγε την απάντηση, όπως \code{bfs_depth_5}, \code{ida_star_depth_10} ή \code{combined_table_lower_bound}, ώστε για κάθε τιμή απόστασης να φαίνεται με ποια μέθοδο προέκυψε.

\section{Πεδία επιλυτή}
\label{sec:schema-solver}
Το πεδίο \code{solver} δηλώνει ποιος επιλυτής παρήγαγε το αποτέλεσμα. Το \code{solution} είναι συμβολοσειρά κινήσεων και το \code{solution_length} το μήκος της αφού αναλυθεί η συμβολοσειρά, στη μετρική μισών στροφών (\eng{half-turn metric}, HTM). Το \code{metric} είναι HTM για όλες τις γραμμές αναφοράς της εργασίας. Τα πεδία \code{runtime_seconds}, \code{expanded_nodes}, \code{generated_nodes} και \code{table_size_bytes} είναι μετρήσεις εκτέλεσης· μπορούν να είναι κενά όταν ο επιλυτής δεν παρέχει την αντίστοιχη πληροφορία.

\section{Πεδία εγγύησης}
\label{sec:schema-status}
Τα πεδία \code{status} και \code{verified} πρέπει να διαβάζονται μαζί. Το \code{verified=true} σημαίνει ότι η λύση εφαρμόστηκε επιτυχώς στην αρχική κατάσταση από τον ανεξάρτητο επαληθευτή. Το \code{status} δείχνει πόσο ισχυρό είναι το αποτέλεσμα. Ο πλήρης κατάλογος των ετικετών κατάστασης, στον οποίο παραπέμπουν τα κεφάλαια του κυρίως σώματος, είναι ο ακόλουθος:

\begin{itemize}
\item \code{exact}: η λύση συνοδεύεται από απόδειξη ότι έχει το ελάχιστο δυνατό μήκος, εντός των ορίων της αναζήτησης· είναι αποδεδειγμένα βέλτιστη.
\item \code{non_exact}: έγκυρη, επαληθευμένη λύση χωρίς απόδειξη βελτιστότητας.
\item \code{lower_bound}: αποδεδειγμένο κάτω φράγμα απόστασης χωρίς λύση στη γραμμή.
\item \code{timeout}: η αναζήτηση σταμάτησε από όριο χρόνου, κόμβων ή βάθους· δεν υπάρχει λύση στη γραμμή, χωρίς αυτό να σημαίνει ότι η κατάσταση είναι άλυτη.
\item \code{not_applicable}: ο επιλυτής δεν υποστηρίζει τη συγκεκριμένη περίπτωση ή απουσιάζει προαιρετική εξάρτηση.
\item \code{failed}: πραγματική αποτυχία, όπως ακολουθία που δεν επαληθεύεται.
\end{itemize}

Η διάκριση ανάμεσα στο \code{timeout} και στο \code{lower_bound} είναι κρίσιμη: η πρώτη ετικέτα δηλώνει απλώς ότι εξαντλήθηκαν οι διαθέσιμοι υπολογιστικοί πόροι (χρόνος ή κόμβοι), ενώ η δεύτερη δίνει αποδεδειγμένη πληροφορία για την απόσταση.

\section{Το πεδίο ελεύθερου κειμένου}
\label{sec:schema-notes}
Το πεδίο \code{notes} είναι ελεύθερο κείμενο με τεχνικά σχόλια. Χρησιμοποιείται για να καταγράψει μήκη φάσεων, την αιτία ενός \code{timeout} ή την απουσία προαιρετικού πακέτου. Δεν πρέπει να χρησιμοποιείται ως μοναδική πηγή δομημένης πληροφορίας: αν θέλουμε να φιλτράρουμε ή να συνοψίσουμε μια πληροφορία σε πίνακα, πρέπει να γίνει ξεχωριστό, δομημένο πεδίο σε μελλοντική έκδοση του σχήματος.
